Negli ultimi anni, è sorto un crescente interesse verso i veicoli aerei a guida autonoma (\textit{Unmanned Aerial Vehicles}, UAV) grazie all'ampio ventaglio di applicazioni sia in ambito civile che militare. 
Trovano infatti applicazione in operazioni di ricerca e soccorso, sorveglianza, ma anche in settori come l'agricoltura e il monitoraggio ambientale.\\
Il principale vantaggio consiste nella loro capacità di muoversi in ambienti complessi e dinamici in maniera autonoma, anche in situazioni che per l'uomo potrebbero essere ostiche o pericolose.\\
Gli UAV sono generalmente classificati in due categorie:
\begin{itemize}
    \item \textbf{Multirotori}: sono caratterizzati da alta stabilità in hovering e capacità di decollare e atterrare verticalmente, ma non sono in grado di percorrere lunghe distanze;
    \item \textbf{Aeromobili ad ala fissa}: presentano alta efficienza aerodinamica e sono in grado di effettuare lunghe percorrenze, ma hanno necessità di infrastrutture dedicate per il decollo e l'atterraggio.
\end{itemize}
Le limitazioni di entrambe le configurazioni vengono superate dai VTOL (Vertical Take-Off and Landing), ovvero UAV a configurazione ibrida, i quali presentano superfici aerodinamiche e rotori inclinabili;
in questo modo offrono la versatilità necessaria per operare in ambienti confinati senza rinunciare a lunghe autonomie di volo.

\section{Il Problema Ingegneristico}
La progettazione di sistemi di controllo per UAV VTOL ibridi presenta sfide non banali, derivanti principalmente dalla natura fortemente non lineare e tempo-variante della loro dinamica.\\
Durante la transizione tra il volo verticale (hovering) e quello orizzontale (crociera), il velivolo attraversa regimi aerodinamici complessi in cui le forze di portanza, resistenza e i momenti giroscopici risultano strettamente accoppiati.\\
Il velivolo oggetto di questo studio è un VTOL innovativo, caratterizzato dalla presenza di ali e tre rotori inclinabili: due anteriori e uno posteriore.
Ciò che rende particolarmente sfidante il controllo dell'inviluppo di volo di questo drone è l'instabilità intrinseca che deriva dall'ultilizzo di tre soli rotori e l'assenza di superfici aerodinamiche a geometria varibile, le quali sono generalmente usate 
per il controllo di assetto in fase di crociera.
\begin{figure}
    \centering
    \includegraphics[width=0.8\textwidth]{dragonfly.png}
    \caption{Dragonfly}
\end{figure}  

\section{Obiettivo della Tesi}
Il presente lavoro si propone di modificare il modello matematico completo a sei gradi di libertà (6-DOF) basato sulla formulazione di Newton-Eulero e di sintetizzare un'architettura di controllo robusta capace di gestire l'intero inviluppo di volo.
\\La strategia proposta si articola su un sistema di controllo gerarchico a cascata:
\begin{enumerate}
    \item Un \textit{Outer Loop} per la guida del drone, con scelte di controllori differenti tra la fase di hovering e quella di crociera.
    \item Un \textit{Inner Loop} di assetto per la stabilizzazione dei momenti.
    \item Un algoritmo di \textit{Motor Mixing} sequenziale che traduce i comandi virtuali in riferimenti fisici per i rotori e i servomotori, rispettando i vincoli di equilibrio meccanico del sistema.
\end{enumerate}

\section{Architettura di Controllo}
L'architettura non convenzionale del velivolo, caratterizzata da una configurazione a tricottero con rotori anteriori avanzati e un rotore di coda a due gradi di libertà, richiede una gestione dinamica della stabilità differenziata per le diverse condizioni operative.

\subsection{Fase 1: Volo Verticale (Hovering)}
In questa configurazione, la portanza è generata interamente dai tre rotori. La sfida principale risiede nel \textbf{bilanciamento del momento di imbardata}. Poiché i rotori anteriori sono controrotanti, la coppia di reazione residua prodotta dal motore posteriore deve essere compensata attivamente per evitare rotazioni indesiderate attorno all'asse $Z_b$.
\begin{itemize}
    \item \textbf{Strategia:} Si utilizza il grado di libertà di tilting laterale del rotore di coda ($\theta_3$) per generare una componente di spinta trasversale. Tale forza crea un momento di imbardata uguale e opposto alla coppia resistente del motore, garantendo l'equilibrio statico dell'assetto.
\end{itemize}

\subsection{Fase 2: Transizione}
Rappresenta la fase a maggiore criticità dinamica, in cui i rotori anteriori ruotano progressivamente verso la configurazione orizzontale per generare spinta propulsiva.
\begin{itemize}
    \item \textbf{Sfide:} Il sistema è soggetto a forti non linearità dovute all'interferenza aerodinamica dei flussi dei rotori sulle superfici alari e alla variazione del braccio di leva longitudinale del vettore spinta rispetto al baricentro.
    \item \textbf{Strategia:} Il drone raggiunge una quota maggiore di quella desiderata, in modo da compensare eventuali perdite durante il \textit{tilting} dei rotori. I rotori anteriori vengono inclinati in avanti e il rotore di coda viene progressivamente spento, fino allo spegnimento definitivo quando il velivolo ha raggiunto la velocità di crociera con cui si genera portanza sufficiente a sostenere il peso del drone.
\end{itemize}

\subsection{Fase 3: Volo Orizzontale (Crociera)}
Al raggiungimento della velocità di sostentamento ($v > v_{stall}$), l'ala assume il ruolo di superficie portante principale. 
\begin{itemize}
    \item \textbf{Strategia:} Il rotore posteriore viene disattivato ($\omega_3 = 0$). L'autorità di controllo è riallocata esclusivamente sui due rotori anteriori:
    \begin{itemize}
        \item \textbf{Rollio:} Gestito tramite tilting differenziale ($\theta_1 \neq \theta_2$).
        \item \textbf{Imbardata:} Gestita tramite modulazione differenziale della spinta ($\Delta \omega_{1,2}$).
        \item \textbf{Beccheggio:} Affidato alla stabilità intrinseca del profilo aerodinamico e a correzioni della spinta propulsiva.
    \end{itemize}
\end{itemize}

Di seguito si trattano dettagliatamente i controllori delle fasi di volo verticale e orizzontale; si lascia lo sviluppo del controllo di trasizione ad eventuali progetti futuri.
